\documentclass[12pt]{ximera}
\typeout{Start loading xmPreamble.tex}%

% Add here extra macro's that are loaded automatically by all documents of claas 'ximera' or 'xourse' in this repo

%%
%%  Example:
%%
% \newcommand{\R}{\mathbb{R}}
\usepackage{tikz}
\usetikzlibrary{positioning, arrows.meta}
\usepackage{multicol}
\usepackage{enumitem}
\usepackage{caption}
\usepackage{amsmath}
\usepackage{amsthm}
\usepackage{fontenc}

%% ---------------------------------------------------------------
%% Shared worksheet macros.
%%
%% These came from the Summer 2026 worksheet set, where each worksheet carried its
%% own copy. They have to live here instead: a xourse inlines every activity into a
%% single document, so duplicate \newcommand definitions across activities collide,
%% and the xourse gobbles \input so a per-topic macro file never loads.
%%
%% They use \providecommand, NOT \newcommand. ximera.cls loads this file automatically,
%% and every activity also carries an explicit \typeout{Start loading xmPreamble.tex}%

% Add here extra macro's that are loaded automatically by all documents of claas 'ximera' or 'xourse' in this repo

%%
%%  Example:
%%
% \newcommand{\R}{\mathbb{R}}
\usepackage{tikz}
\usetikzlibrary{positioning, arrows.meta}
\usepackage{multicol}
\usepackage{enumitem}
\usepackage{caption}
\usepackage{amsmath}
\usepackage{amsthm}
\usepackage{fontenc}

%% ---------------------------------------------------------------
%% Shared worksheet macros.
%%
%% These came from the Summer 2026 worksheet set, where each worksheet carried its
%% own copy. They have to live here instead: a xourse inlines every activity into a
%% single document, so duplicate \newcommand definitions across activities collide,
%% and the xourse gobbles \input so a per-topic macro file never loads.
%%
%% They use \providecommand, NOT \newcommand. ximera.cls loads this file automatically,
%% and every activity also carries an explicit \typeout{Start loading xmPreamble.tex}%

% Add here extra macro's that are loaded automatically by all documents of claas 'ximera' or 'xourse' in this repo

%%
%%  Example:
%%
% \newcommand{\R}{\mathbb{R}}
\usepackage{tikz}
\usetikzlibrary{positioning, arrows.meta}
\usepackage{multicol}
\usepackage{enumitem}
\usepackage{caption}
\usepackage{amsmath}
\usepackage{amsthm}
\usepackage{fontenc}

%% ---------------------------------------------------------------
%% Shared worksheet macros.
%%
%% These came from the Summer 2026 worksheet set, where each worksheet carried its
%% own copy. They have to live here instead: a xourse inlines every activity into a
%% single document, so duplicate \newcommand definitions across activities collide,
%% and the xourse gobbles \input so a per-topic macro file never loads.
%%
%% They use \providecommand, NOT \newcommand. ximera.cls loads this file automatically,
%% and every activity also carries an explicit \input{../xmPreamble}, so this file is
%% read twice. That was harmless while it held only \usepackage lines (LaTeX skips a
%% package it has already loaded) but a second \newcommand is a hard error.
%%
%%   \blank[width]        fill-in-the-blank rule (default 2.5cm)
%%   \showwork            "show and explain your work" reminder
%%   \showworkline        ... plus which schedule line was used
%%   \showworkyear        ... plus which year's schedule was used
%%   \schedhead{...}      centred bold caption above a tiered-rate table
%%   \samesquares[scale]{left}{right}   two equal-area squares, labelled
%%   \samecubes[scale]{left}{right}     two equal-volume cubes, labelled
%% ---------------------------------------------------------------

\providecommand{\blank}[1][2.5cm]{\underline{\hspace{#1}}}
\providecommand{\showwork}{\textbf{REMEMBER TO SHOW AND EXPLAIN YOUR WORK.}}
\providecommand{\showworkline}{\textbf{REMEMBER TO SHOW AND EXPLAIN YOUR WORK.
  Explanation should include how and why you know which line of the
  schedule was used to calculate this amount.}}
\providecommand{\showworkyear}{\begin{sloppypar}\textbf{REMEMBER TO SHOW AND
  EXPLAIN YOUR WORK. Explanation should include how and why you know which
  line of the year's tax schedule was used to calculate this tax
  amount.}\end{sloppypar}}
\providecommand{\schedhead}[1]{\begin{center}\textbf{#1}\end{center}}

\providecommand{\samesquares}[3][1]{%
\begin{center}
{\footnotesize These squares are the \textbf{same size}.}

\vspace{1.5mm}
\begin{tikzpicture}[scale=#1,every node/.style={scale=#1}]
  \draw (0,0) rectangle (2.4,2.4);
  \node[anchor=north west] at (0.15,2.25) {Area:};
  \node[anchor=east]  at (-0.15,1.2) {#2};
  \node[anchor=north] at (1.2,-0.15) {#2};
  \begin{scope}[xshift=4.6cm]
    \draw (0,0) rectangle (2.4,2.4);
    \node[anchor=north west] at (0.15,2.25) {Area:};
    \node[anchor=east]  at (-0.15,1.2) {#3};
    \node[anchor=north] at (1.2,-0.15) {#3};
  \end{scope}
\end{tikzpicture}
\end{center}}

\providecommand{\samecubes}[3][1]{%
\begin{center}
{\footnotesize These cubes are the \textbf{same size}.}

\vspace{1.5mm}
\begin{tikzpicture}[scale=#1,every node/.style={scale=#1}]
  \foreach \dx in {0,5.2} {
    \begin{scope}[xshift=\dx cm]
      \draw (0,0) rectangle (2.0,2.0);
      \draw (0,2.0) -- (0.65,2.6) -- (2.65,2.6) -- (2.0,2.0);
      \fill[black!12] (2.0,0) -- (2.65,0.6) -- (2.65,2.6) -- (2.0,2.0) -- cycle;
      \draw (2.0,0) -- (2.65,0.6) -- (2.65,2.6) -- (2.0,2.0);
      \node[anchor=north west] at (0.1,1.9) {Volume:};
    \end{scope}
  }
  \node[anchor=east]  at (-0.15,1.0) {#2};
  \node[anchor=north] at (1.0,-0.15) {#2};
  \node[anchor=west]  at (2.25,0.4) {#2};
  \node[anchor=east]  at (5.05,1.0) {#3};
  \node[anchor=north] at (6.2,-0.15) {#3};
  \node[anchor=west]  at (7.45,0.4) {#3};
\end{tikzpicture}
\end{center}}

%% NOTE: do NOT add \usepackage to individual activity .tex files.
%% When a xourse inlines an activity it gobbles \documentclass and \input, but a bare
%% \usepackage still executes -- after \begin{document} -- and errors out.
%% All shared packages and macros belong here.
%%
%% ximera.cls already loads: amsmath, amssymb, amsthm, cancel, enumitem, graphicx,
%% hyperref, listings, pgfplots, tikz, url, xcolor. No need to re-load those.
, so this file is
%% read twice. That was harmless while it held only \usepackage lines (LaTeX skips a
%% package it has already loaded) but a second \newcommand is a hard error.
%%
%%   \blank[width]        fill-in-the-blank rule (default 2.5cm)
%%   \showwork            "show and explain your work" reminder
%%   \showworkline        ... plus which schedule line was used
%%   \showworkyear        ... plus which year's schedule was used
%%   \schedhead{...}      centred bold caption above a tiered-rate table
%%   \samesquares[scale]{left}{right}   two equal-area squares, labelled
%%   \samecubes[scale]{left}{right}     two equal-volume cubes, labelled
%% ---------------------------------------------------------------

\providecommand{\blank}[1][2.5cm]{\underline{\hspace{#1}}}
\providecommand{\showwork}{\textbf{REMEMBER TO SHOW AND EXPLAIN YOUR WORK.}}
\providecommand{\showworkline}{\textbf{REMEMBER TO SHOW AND EXPLAIN YOUR WORK.
  Explanation should include how and why you know which line of the
  schedule was used to calculate this amount.}}
\providecommand{\showworkyear}{\begin{sloppypar}\textbf{REMEMBER TO SHOW AND
  EXPLAIN YOUR WORK. Explanation should include how and why you know which
  line of the year's tax schedule was used to calculate this tax
  amount.}\end{sloppypar}}
\providecommand{\schedhead}[1]{\begin{center}\textbf{#1}\end{center}}

\providecommand{\samesquares}[3][1]{%
\begin{center}
{\footnotesize These squares are the \textbf{same size}.}

\vspace{1.5mm}
\begin{tikzpicture}[scale=#1,every node/.style={scale=#1}]
  \draw (0,0) rectangle (2.4,2.4);
  \node[anchor=north west] at (0.15,2.25) {Area:};
  \node[anchor=east]  at (-0.15,1.2) {#2};
  \node[anchor=north] at (1.2,-0.15) {#2};
  \begin{scope}[xshift=4.6cm]
    \draw (0,0) rectangle (2.4,2.4);
    \node[anchor=north west] at (0.15,2.25) {Area:};
    \node[anchor=east]  at (-0.15,1.2) {#3};
    \node[anchor=north] at (1.2,-0.15) {#3};
  \end{scope}
\end{tikzpicture}
\end{center}}

\providecommand{\samecubes}[3][1]{%
\begin{center}
{\footnotesize These cubes are the \textbf{same size}.}

\vspace{1.5mm}
\begin{tikzpicture}[scale=#1,every node/.style={scale=#1}]
  \foreach \dx in {0,5.2} {
    \begin{scope}[xshift=\dx cm]
      \draw (0,0) rectangle (2.0,2.0);
      \draw (0,2.0) -- (0.65,2.6) -- (2.65,2.6) -- (2.0,2.0);
      \fill[black!12] (2.0,0) -- (2.65,0.6) -- (2.65,2.6) -- (2.0,2.0) -- cycle;
      \draw (2.0,0) -- (2.65,0.6) -- (2.65,2.6) -- (2.0,2.0);
      \node[anchor=north west] at (0.1,1.9) {Volume:};
    \end{scope}
  }
  \node[anchor=east]  at (-0.15,1.0) {#2};
  \node[anchor=north] at (1.0,-0.15) {#2};
  \node[anchor=west]  at (2.25,0.4) {#2};
  \node[anchor=east]  at (5.05,1.0) {#3};
  \node[anchor=north] at (6.2,-0.15) {#3};
  \node[anchor=west]  at (7.45,0.4) {#3};
\end{tikzpicture}
\end{center}}

%% NOTE: do NOT add \usepackage to individual activity .tex files.
%% When a xourse inlines an activity it gobbles \documentclass and \input, but a bare
%% \usepackage still executes -- after \begin{document} -- and errors out.
%% All shared packages and macros belong here.
%%
%% ximera.cls already loads: amsmath, amssymb, amsthm, cancel, enumitem, graphicx,
%% hyperref, listings, pgfplots, tikz, url, xcolor. No need to re-load those.
, so this file is
%% read twice. That was harmless while it held only \usepackage lines (LaTeX skips a
%% package it has already loaded) but a second \newcommand is a hard error.
%%
%%   \blank[width]        fill-in-the-blank rule (default 2.5cm)
%%   \showwork            "show and explain your work" reminder
%%   \showworkline        ... plus which schedule line was used
%%   \showworkyear        ... plus which year's schedule was used
%%   \schedhead{...}      centred bold caption above a tiered-rate table
%%   \samesquares[scale]{left}{right}   two equal-area squares, labelled
%%   \samecubes[scale]{left}{right}     two equal-volume cubes, labelled
%% ---------------------------------------------------------------

\providecommand{\blank}[1][2.5cm]{\underline{\hspace{#1}}}
\providecommand{\showwork}{\textbf{REMEMBER TO SHOW AND EXPLAIN YOUR WORK.}}
\providecommand{\showworkline}{\textbf{REMEMBER TO SHOW AND EXPLAIN YOUR WORK.
  Explanation should include how and why you know which line of the
  schedule was used to calculate this amount.}}
\providecommand{\showworkyear}{\begin{sloppypar}\textbf{REMEMBER TO SHOW AND
  EXPLAIN YOUR WORK. Explanation should include how and why you know which
  line of the year's tax schedule was used to calculate this tax
  amount.}\end{sloppypar}}
\providecommand{\schedhead}[1]{\begin{center}\textbf{#1}\end{center}}

\providecommand{\samesquares}[3][1]{%
\begin{center}
{\footnotesize These squares are the \textbf{same size}.}

\vspace{1.5mm}
\begin{tikzpicture}[scale=#1,every node/.style={scale=#1}]
  \draw (0,0) rectangle (2.4,2.4);
  \node[anchor=north west] at (0.15,2.25) {Area:};
  \node[anchor=east]  at (-0.15,1.2) {#2};
  \node[anchor=north] at (1.2,-0.15) {#2};
  \begin{scope}[xshift=4.6cm]
    \draw (0,0) rectangle (2.4,2.4);
    \node[anchor=north west] at (0.15,2.25) {Area:};
    \node[anchor=east]  at (-0.15,1.2) {#3};
    \node[anchor=north] at (1.2,-0.15) {#3};
  \end{scope}
\end{tikzpicture}
\end{center}}

\providecommand{\samecubes}[3][1]{%
\begin{center}
{\footnotesize These cubes are the \textbf{same size}.}

\vspace{1.5mm}
\begin{tikzpicture}[scale=#1,every node/.style={scale=#1}]
  \foreach \dx in {0,5.2} {
    \begin{scope}[xshift=\dx cm]
      \draw (0,0) rectangle (2.0,2.0);
      \draw (0,2.0) -- (0.65,2.6) -- (2.65,2.6) -- (2.0,2.0);
      \fill[black!12] (2.0,0) -- (2.65,0.6) -- (2.65,2.6) -- (2.0,2.0) -- cycle;
      \draw (2.0,0) -- (2.65,0.6) -- (2.65,2.6) -- (2.0,2.0);
      \node[anchor=north west] at (0.1,1.9) {Volume:};
    \end{scope}
  }
  \node[anchor=east]  at (-0.15,1.0) {#2};
  \node[anchor=north] at (1.0,-0.15) {#2};
  \node[anchor=west]  at (2.25,0.4) {#2};
  \node[anchor=east]  at (5.05,1.0) {#3};
  \node[anchor=north] at (6.2,-0.15) {#3};
  \node[anchor=west]  at (7.45,0.4) {#3};
\end{tikzpicture}
\end{center}}

%% NOTE: do NOT add \usepackage to individual activity .tex files.
%% When a xourse inlines an activity it gobbles \documentclass and \input, but a bare
%% \usepackage still executes -- after \begin{document} -- and errors out.
%% All shared packages and macros belong here.
%%
%% ximera.cls already loads: amsmath, amssymb, amsthm, cancel, enumitem, graphicx,
%% hyperref, listings, pgfplots, tikz, url, xcolor. No need to re-load those.

\title{In-Class: Exponentials, Part 1}
\begin{document}
\begin{abstract}
Building an exponential function $y = a \cdot R^x$ from a repeated percent increase,
interpreting each piece in context, and using Guess \& Check to estimate doubling
times for a CD investment and an epidemic.
\end{abstract}
\maketitle

\section*{Where the exponential comes from}

We are going to revisit some scenarios from the pre-class work and try to understand
how they are ``exponential,'' and how we can write and interpret an exponential
function for those scenarios.

\begin{example}
An athlete signs a contract saying they will earn \$8.3 million with an increase of
$4.8\%$ each year of the 5 year contract.

\[
\begin{aligned}
\text{year 0:} \quad & \$8.3 \text{ mil} \\
\text{year 1:} \quad & \$8.3 \text{ mil} + 0.048 \cdot (\$8.3 \text{ mil})
  = \$8.3 \text{ mil} \cdot (1 + 0.048) = \$8.3 \text{ mil} \cdot (1.048) \\
\text{year 2:} \quad & \$8.3 \text{ mil} \cdot (1.048) + 0.048 \cdot \left(\$8.3 \text{ mil} \cdot (1.048)\right) \\
  &= \$8.3 \text{ mil} \cdot (1.048)\left[1 + 0.048\right] = \$8.3 \text{ mil} \cdot (1.048)^2 \\
\text{year 3:} \quad & \$8.3 \text{ mil} \cdot (1.048)^2 \left[1 + 0.048\right] = \$8.3 \text{ mil} \cdot (1.048)^3 \\
\text{year 4:} \quad & \$8.3 \text{ mil} \cdot (1.048)^3 \left[1 + 0.048\right] = \$8.3 \text{ mil} \cdot (1.048)^4
\end{aligned}
\]

Each year we multiply by the \emph{same} factor $(1 + 0.048)$, so after $x$ years we
have multiplied by it $x$ times. That repeated multiplication is what makes the
situation exponential.
\end{example}

\begin{problem}
You filled out the first two columns of this table in your pre-class work. From the
calculations above, fill in the last column by noticing the pattern.

\begin{center}
\begin{tabular}{|c|l|l|l|}
\hline
After \_\_ years & Salary & From the previous year's salary & From the starting salary \\
\hline
0 & \$8.3 mill & N/A & $8.3(1.048)^{\answer{0}}$ \\
\hline
1 & \$8.7 mill & $\$8.3(0.048) + \$8.3 = \$8.3(1+0.048)$ & $8.3(1.048)^{\answer{1}}$ \\
\hline
2 & \$9.1 mill & $\$8.7(1+0.048)$ & $8.3(1.048)^{\answer{2}}$ \\
\hline
3 & \$9.5 mill & $\$9.1(1+0.048)$ & $8.3(1.048)^{\answer{3}}$ \\
\hline
4 & \$10.0 mill & $\$9.5(1+0.048)$ & $8.3(1.048)^{\answer{4}}$ \\
\hline
\end{tabular}
\end{center}
\end{problem}

\begin{problem}
Suppose the athlete renewed this contract indefinitely, so the $4.8\%$ yearly
increase didn't stop after 5 years.

\begin{prompt}
\textbf{(a)} How much would the athlete earn at the end of the 17th year? Give your
answer in millions of dollars.
\[
\$\answer[tolerance=0.05]{18.42} \text{ million}
\]
\begin{solution}
The salary would be $\$8.3 \text{ mil} \cdot (1.048)^{17} \approx \$18.42$ million.
\end{solution}
\end{prompt}

\begin{prompt}
\textbf{(b)} How much would the athlete be earning after the $x$th year?
\[
y = \answer{8.3}(\answer{1.048})^{x}
\]
\end{prompt}
\end{problem}

\subsection*{Reading the function}

\emph{Do not read further if you have not completed the work above --- ask your
instructor if you need help.}

What we developed above is an exponential function describing the athlete's salary.
Let's look at each value and identify what it represents.

\[
y = 8.3(1.048)^{x}
\]

\begin{description}
\item[$8.3$] \emph{Units:} millions of dollars.

  \emph{Represents:} the $y$-value (output) when the $x$-value (input) is $0$ ---
  the salary when the number of years since signing is $0$. Note that when $x = 0$,
  $y = \$8.3 \text{ mil} \cdot (1.048)^0$, and $(1.048)^0 = 1$, so
  $y = \$8.3$ mil. So \$8.3 million is the \textbf{initial salary} of the athlete.

\item[$1.048$] \emph{Rate of change.} No units, since $1.048$ comes from a percent.

  \emph{Represents:} the athlete's salary increases by $4.8\%$ each year.

\item[$x$] \emph{Units:} years. The number of years since the start of the contract.

\item[$y$] \emph{Units:} millions of dollars. The athlete's salary after $x$ years.
\end{description}

This is the general form $y = a \cdot R^{x}$. Now we can write exponential functions
for new scenarios without going through the step-by-step table each time.

\section*{The CD investment}

Some CDs (a banking investment option) offer rates giving about $3.10\%$ yield per
year, as long as a minimum amount such as \$5,000 is invested. ``$3.10\%$ yield''
simply means the investment earns about $3.10\%$. CDs run for a certain number of
years; when a CD ``matures'' you usually have the option to renew it.

\begin{problem}
Suppose you invest the minimum \$5,000 to get this $3.10\%$ yield. Write a function
describing what the investment is worth after $x$ years.
\[
y = \answer{5000}(\answer{1.031})^{x}
\]

\begin{solution}
\$5000 is the $y$-value when $x = 0$, since $y = \$5000 \cdot (1.031)^0 = \$5000$ ---
the value of the CD when 0 years have passed, i.e.\ the initial amount invested.

$1.031$ tells us the value of the CD increases by $3.10\%$ each year.

$x$ is the number of years since the start of the CD, and $y$ is the value of the CD
in dollars after $x$ years.
\end{solution}
\end{problem}

\begin{problem}
Fill in the first column with the values from your pre-class work, then use your
function to fill in the last column.

\begin{center}
\begin{tabular}{|c|l|l|}
\hline
After year & Value from Pre-Class & Value from the function \\
\hline
0 & \$5,000 & $\$\answer{5000}$ \\
\hline
1 & \$5,155 & $\$\answer[tolerance=0.5]{5155}$ \\
\hline
2 & \$5,314.81 & $\$\answer[tolerance=0.5]{5314.81}$ \\
\hline
3 & \$5,479.57 & $\$\answer[tolerance=0.5]{5479.57}$ \\
\hline
\end{tabular}
\end{center}

How do these values compare? Are they exactly the same? Should they be? Explain.

\begin{freeResponse}
\end{freeResponse}
\end{problem}

\begin{problem}
Suppose this is a 7 year CD. How much is the investment worth at the end?
\[
\$\answer[tolerance=0.5]{6191.28}
\]
\begin{solution}
When $x = 7$, $y = \$5000 \cdot (1.031)^{7} = \$6{,}191.28$.
\end{solution}
\end{problem}

\subsection*{Guess \& Check}

\begin{problem}
Suppose you keep renewing this CD at this rate every time it matures. About how many
years will it take to double the initial investment? Use \textbf{Guess and Check}.

It will take about $\answer{23}$ years.

\begin{hint}
The initial investment was \$5000, so doubling it means the CD must reach \$10,000.
We want the $x$ satisfying $\$5000 \cdot (1.031)^{x} = \$10{,}000$.
\end{hint}

\begin{solution}
\textbf{Step 1: find the closest whole number giving \emph{less} than we want.}
\[
x = 22: \quad \$5000 \cdot (1.031)^{22} = \$9{,}787.25
\]

\textbf{Step 2: find the closest whole number giving \emph{more} than we want.}
\[
x = 23: \quad \$5000 \cdot (1.031)^{23} = \$10{,}090.65
\]

\textbf{Step 3: which is closer to what we want?}
\[
\begin{array}{ll}
x = 22 & \$10{,}000 - \$9{,}787.25 = \$212.75 \\
x = 23 & \$10{,}090.65 - \$10{,}000 = \$90.65
\end{array}
\]
\$10,090.65 is closer to the desired \$10,000, so we use $23$ as our estimate. It
will take approximately 23 years for the initial investment to double.
\end{solution}
\end{problem}

This problem outlined the ``steps'' for carrying out Guess \& Check. In future
problems these steps will NOT be listed for you --- but they show how much work you
need to show to receive full credit for a Guess \& Check problem.

\section*{An epidemic model}

Using CDC data about ebola cases in Guinea, an exponential model can be created
representing the number of cases $x$ days after May 25th, 2014, through the remainder
of 2014:
\[
y = 121.2(1.0117)^{x}
\]

\begin{problem}
\begin{prompt}
\textbf{(a)} What does $x$ represent in this context?
\begin{freeResponse}
\end{freeResponse}
\begin{solution}
$x$ is the number of days that have passed since May 25, 2014.
\end{solution}
\end{prompt}

\begin{prompt}
\textbf{(b)} What does $y$ represent?
\begin{freeResponse}
\end{freeResponse}
\begin{solution}
$y$ is the number of ebola cases $x$ days after May 25, 2014.
\end{solution}
\end{prompt}

\begin{prompt}
\textbf{(c)} What does $121.2$ represent?
\begin{freeResponse}
\end{freeResponse}
\begin{solution}
$121.2$ is the $y$-value when $x = 0$, so it is the number of ebola cases on
May 25, 2014.
\end{solution}
\end{prompt}

\begin{prompt}
\textbf{(d)} What does $1.0117$ tell us about this context? The number of cases
increases by $\answer[tolerance=0.005]{1.17}$\% each day.
\end{prompt}
\end{problem}

\begin{problem}
Use Guess \& Check to estimate when the number of ebola cases doubles. Show the same
steps outlined above.

The number of cases doubles after about $\answer{60}$ days.

\begin{hint}
There were initially $121.2$ cases, so doubling gives $242.4$ cases. We want the $x$
satisfying $121.2(1.0117)^{x} = 242.4$.
\end{hint}

\begin{solution}
\textbf{Step 1.} Choose the \emph{largest} whole number $x$ for which the output is
\emph{less} than $242.4$:
\[
x = 59: \quad 121.2 \cdot (1.0117)^{59} = 240.7
\]

\textbf{Step 2.} Choose the \emph{smallest} whole number $x$ for which the output is
\emph{more} than $242.4$:
\[
x = 60: \quad 121.2 \cdot (1.0117)^{60} = 243.6
\]

\textbf{Step 3.} Which is closer?
\[
\begin{aligned}
x = 59: \quad & 242.4 - 240.7 = 1.7 \\
x = 60: \quad & 243.6 - 242.4 = 1.2
\end{aligned}
\]
$243.6$ is closer to $242.4$, so we take \textbf{60 days} as the estimate for the
number of days after May 25, 2014 for the cases to double.
\end{solution}
\end{problem}

\begin{problem}
Suppose there is a new virus that reportedly doubles infection cases in about 20
days. What would this virus' infection rate be, as a percent?

The infection rate is about $\answer[tolerance=0.1]{3.5}$\% per day.

\begin{hint}
Set up an exponential function assuming there were initially 80 recorded infections.
\end{hint}

\begin{solution}
We do not know the daily rate; call it $r$. With 80 initial cases doubling to 160
after 20 days,
\[
80(1+r)^{20} = 160.
\]
Then $(1+r)^{20} = 2$, and taking the 20th root of both sides,
\[
1 + r = 2^{1/20}, \qquad r = 2^{1/20} - 1 \approx 1.035 - 1 = 0.035.
\]
So the number of infections increases by about $3.5\%$ each day.
\end{solution}
\end{problem}

\begin{problem}
Why was the initial number of infections not actually needed to answer the previous
question?

\begin{freeResponse}
\end{freeResponse}

\begin{solution}
The 80 divides out. Starting from $a(1+r)^{20} = 2a$, dividing both sides by $a$
leaves $(1+r)^{20} = 2$ no matter what $a$ is. Doubling time depends only on the
rate, not on how many cases you start with.
\end{solution}
\end{problem}

\end{document}
